\documentclass[a4paper,10pt]{article}
\usepackage[francais]{babel}
\usepackage{url}

\selectlanguage{francais}
\input{lettre-ensiie.sty}

\selectlanguage{francais}
\enCouleur

\begin{document}


\name{St\'ephane \textsc{Genaud}}
\fonction{Professeur en Informatique à l'ENSIIE\\Directeur du Campus de Strasbourg}
\telephoneillkirch{+33 (0)369 200 231}
\email{genaud@ensiie.fr}

\begin{letter}[Avis concernant l'inscription à l'HDR de Karine Deschinkel]{à}{Service Recherche de l'Université de Franche-Comté}

Madame, Monsieur,

Après une formation supérieure en mathématiques appliquées,
Karine Deschinkel a obtenu un doctorat en système industriels en 2001 de 
l'Ecole Nationale Supérieure de l'Aéronautique et de l'Espace (ENSAE) à Toulouse.
Elle a ensuite occupé un poste d'ATER (section CNU 27) à Paris 1, puis a été recrutée
comme Maître de conférences à l'Université de Versailles Saint Quentin en Yvelines
de 2002 à 2009. Enfin, elle a rejoint l'Université de Franche-Comté sur un poste de 
Maître de conférences en 2009.

Ce parcours riche en expériences professionnelles dans des environnements différents,
à la fois géographiques et thématiques, est tout à fait intéressant. Madame Deschnikel
a pu apporter à différents domaines ses compétences en techniques d'optimisation.
Pour les équipes de recherche en informatique, une telle compétence est précieuse.
Dans le poste qu'elle occupe actuellement à l'IUT de Belfort-Montbelliard, les lettres 
de recommandation témoignent d'une très bonne intégration. Les publications nombreuses
au sein de l'équipe montrent également que son travail s'insère parfaitement, dans l'équipe
mais également dans des domaines transverses couverts par FEMTO-ST.
Ses perspectives de recherche au sein de l'équipe AND, sur le thème de l'optimisation de 
la consommation énergétique dans les réseaux de capteurs, sont très pertinentes d'un point 
de vue scientifique. 


Pour toutes ces raisons, j'ai un avis extrêmement favorable à l'inscription de Madame  
Deschinkel en diplôme d'HDR.


Veuillez agréer, Madame,Monsieur, l'expression de mes sentiments respectueux.

\end{letter}
\begin{flushright}
\includegraphics[width=.26\textwidth]{signgenaud.jpg}
\end{flushright}



\end{document}
